\section{Introduction}\label{sec:introduction}

All Banach spaces in this paper are real. A map \(F:X\to Y\) is a
\emph{bi-Lipschitz embedding} if there are constants
\(0<c\leq C<\infty\) such that
\[
 c\|x-x'\|_X\leq\|F(x)-F(x')\|_Y
       \leq C\|x-x'\|_X\qquad(x,x'\in X).
\]
It is a \emph{bi-Lipschitz equivalence} when it is also onto \(Y\).
The distinction between these two notions is central here: the result
constructs an onto map between whole Banach spaces, while its linear
obstruction concerns a subspace.

The nonlinear geometry of Banach spaces asks how much linear structure
is retained by maps that control distances but need not respect
addition. Ribe proved that uniformly homeomorphic normed spaces have
the same finite-dimensional linear structure up to one common
distortion bound \cite[Theorem~1]{Ribe76}. Aharoni and Lindenstrauss
constructed nonseparable Banach spaces that are bi-Lipschitz equivalent
but not linearly isomorphic, and asked for separable examples
\cite[pp.~281--282]{AL78}. Ribe later produced separable uniformly
homeomorphic nonisomorphic spaces, and Aharoni and Lindenstrauss
obtained uniformly convex examples that are uniformly homeomorphic
but not Lipschitz equivalent \cite{Ribe84,AL85}.
Johnson, Lindenstrauss and Schechtman
developed both rigidity results and separable examples with different
linear structures under the weaker relation of uniform homeomorphism
\cite[Theorem~5.8 and Section~8(3)]{JLS96}. These results explain why a separable bi-Lipschitz
counterexample requires more than a uniform deformation.

The space \(c_0=c_0(\N)\) consists of real scalar sequences tending to
zero, with the supremum norm. Godefroy, Kalton and Lancien proved that a
Banach space bi-Lipschitz equivalent to a linear subspace of \(c_0\) is
linearly isomorphic to a subspace of \(c_0\). They also proved that a
Banach space bi-Lipschitz equivalent to \(c_0\) is linearly isomorphic
to \(c_0\) \cite[Theorems~2.1--2.2]{GKL00}. These theorems concern the
whole domain of an equivalence. They do not say that a Banach space
containing a bi-Lipschitz image of \(c_0\) must contain a linear copy
of \(c_0\). Kalton recorded the latter question as Problem~2, separately
from the separable Lipschitz-isomorphism question in Problem~3
\cite[pp.~21--22]{Kalton08}. H\'ajek, Johanis and Schlumprecht still described the
embedding question as apparently open in their January 2024 preprint
(arXiv:2401.00831v1), later published in 2025 \cite{HJS25}.

For Banach spaces \(X,Y\), write \(X\oplus_\infty Y\) for their product
with norm \(\|(x,y)\|=\max\{\|x\|,\|y\|\}\). The main result is the
following whole-space form of a negative answer to the \(c_0\)
embedding question.

\begin{theorem}\label{thm:main}
There is a separable real Banach space \(Z\) containing no closed linear
subspace isomorphic to \(c_0\), and an onto bi-Lipschitz map
\[
                 F:Z\oplus_\infty c_0\longrightarrow Z.
\]
\end{theorem}
